\documentclass[11pt]{article}
\usepackage[utf8]{inputenc}
\usepackage[T1]{fontenc}
\usepackage[brazil]{babel}
\usepackage[margin=2.4cm]{geometry}
\usepackage{amsmath}
\usepackage{booktabs}
\usepackage{parskip}

\title{\textbf{Explicação didática — modelagem do peso de VALE3 nos fundos Itaú}}
\author{TCC — FGV EESP \\ Orientador: Prof.\ Maurício Ferraresi Jr.}
\date{\today}

\begin{document}
\maketitle

Este documento explica, passo a passo e em linguagem acessível, \emph{o que}
estamos fazendo, \emph{por que} cada modelo e cada escolha, e \emph{como}
interpretar os resultados. O companheiro \texttt{results.pdf} traz as tabelas.

\section{O problema}

Queremos entender e, no fim, \textbf{prever} quanto cada fundo da gestora Itaú
investe na ação \textbf{VALE3} — ou seja, o \textbf{peso} de VALE3 na carteira de
cada fundo, mês a mês. A pergunta de fundo é: \emph{quais características dos fundos
explicam essas posições?} Isso se conecta à literatura de \emph{demand-based asset
pricing} (Koijen--Yogo), em que a demanda por um ativo é vista como função de
características de quem investe.

\section{A ideia central: o peso como ``função de demanda''}

Para um fundo $i$ no mês $t$, modelamos:
\[
\text{peso}_{i,t} \;=\; \alpha_t \;+\; \beta^1_t\,x^1_{i,t} \;+\; \dots \;+\;
\beta^k_t\,x^k_{i,t} \;+\; \varepsilon_{i,t},
\]
onde os $x$ são características do fundo (tamanho, nº de cotistas, etc.). Os
coeficientes $\theta_t = (\alpha_t, \beta^1_t, \dots)$ dizem \emph{como} as
características se traduzem em posição naquele mês. A sacada do orientador: em vez
de acompanhar o peso de centenas de fundos, acompanhamos esses poucos coeficientes
— eles são um \textbf{fator latente} que \emph{reduz um problema de alta dimensão a
uma dimensão pequena}.

\section{Por que ``regressão em dois passos''}

\textbf{Passo 1 (transversal / cross-section):} para \emph{cada mês}, uma regressão
separada do peso contra as características dos fundos. Isso produz uma série de
coeficientes $\theta_t$ ao longo do tempo. \\
\textbf{Passo 2 (temporal / dinâmico):} modelar como esses $\theta_t$ \emph{evoluem
no tempo}. O orientador propõe que evoluam como um \emph{passeio aleatório} (random
walk): $\theta_t = \theta_{t-1} + \text{choque}$, estimável pelo \emph{filtro de
Kalman}. É o ``modelo de fator dinâmico''. A previsão vive aqui: prevendo
$\theta_{t+1}$ e conhecendo as características do fundo, reconstruímos o peso
previsto. (Este documento cobre o Passo~1; o Passo~2 é o próximo.)

\section{As variáveis e por que cada tratamento}

\textbf{Alvo: peso (nível).} É uma fração em $[0,1]$. Usamos em nível porque é o que
o orientador pediu (o coeficiente OLS vira o fator) e é direto de interpretar.

\textbf{AUM e nº de cotistas: usamos $\log$.} Essas variáveis são muito
assimétricas — há fundos de R\$ 1 milhão e de R\$ 18 bilhões. Em nível, um fundo
gigante dominaria a regressão. O logaritmo comprime a escala e capta efeitos
\emph{proporcionais}: ir de 10 para 20 milhões ``pesa'' como ir de 100 para 200.

\textbf{FIC vs.\ FI: variável binária (dummy).} Um FIC (Fundo de Investimento em
Cotas) investe em cotas de outros fundos; um FI investe direto. Isso muda a posição
\emph{direta} em ações. Detectamos FIC pelo nome do fundo (a classificação ANBIMA
não distingue FIC/FI).

\textbf{Fluxo líquido: usamos fluxo/AUM.} O fluxo é (captação $-$ resgate) do mês,
em reais. Dividir pelo tamanho do fundo dá o fluxo como \emph{proporção do
patrimônio} — comparável entre fundos grandes e pequenos, e preserva o sinal
(entrada positiva, saída negativa). O fluxo líquido ``de verdade'' (com a perna de
resgate) só existe no \emph{Informe Diário} da CVM, não na base SH.

\section{Passo 1A: a cross-section mês a mês (e por que Fama--MacBeth)}

Rodamos 72 regressões (uma por mês, 2016--2021). Cada uma dá um $\theta_t$. Para
resumir, usamos o procedimento de \textbf{Fama--MacBeth}: a estimativa final de
cada coeficiente é a \emph{média} dos 72 valores mensais, e o erro-padrão vem da
\emph{variação entre os meses} (não de uma única regressão). É o jeito estatístico
correto de tratar coeficientes que mudam mês a mês — e o teste $t$ resultante diz
se o efeito médio é robusto no tempo.

\textbf{Resultado:} tamanho ($-$), cotistas ($+$) e FIC ($-$) têm sinal idêntico
\emph{em todos os 72 meses} e são altamente significativos. O fluxo não é
significativo para o \emph{nível} do peso.

\section{Por que preço e beta NÃO entram na cross-section}

Dentro de um mesmo mês, o preço da VALE3 e o beta da VALE3 são \emph{iguais para
todos os fundos} (são variáveis ``de mercado'', do ativo, não do fundo). Logo, eles
não conseguem explicar por que um fundo difere de outro \emph{naquele mês} — seriam
absorvidos pelo intercepto. São variáveis \emph{agregadas}: explicam como o nível
\emph{médio} das posições se move \emph{ao longo do tempo}. Por isso entram (i) na
regressão \emph{pooled} e (ii) no Passo~2 (a dinâmica do nível médio).

\section{Passo 1B: a regressão \emph{pooled}}

Empilhando todos os fundos-meses numa só regressão, o preço e o beta passam a ter
variação (entre os 72 meses) e ficam identificados. Achamos: \textbf{preço $+$} e
\textbf{beta $-$}, ambos significativos. Ou seja, em meses de VALE3 mais cara, o
peso médio sobe; em meses de VALE3 mais arriscada (beta alto), o peso médio cai
(aversão a risco). \emph{Importante:} no piloto de só 2016 (12 meses), o beta saía
positivo e fraco; com 72 meses ele vira claramente negativo. Isso ensina uma lição
metodológica: 12 pontos no tempo eram poucos para identificar variáveis agregadas
— por isso estendemos para 6 anos.

\section{Como o beta é calculado (e por quê)}

Beta mede \textbf{quanto a ação se move quando o mercado se move}. É a inclinação da
regressão dos retornos da VALE3 sobre os do Ibovespa:
\[
\beta = \frac{\mathrm{Cov}(r_{VALE},\, r_{Ibov})}{\mathrm{Var}(r_{Ibov})}.
\]
Calculamos \emph{nós mesmos} (não pegamos o ``beta'' pronto do Yahoo, que é de 5
anos contra índice americano e estático). Usamos \textbf{retornos diários simples} e
uma \textbf{janela móvel de 252 pregões} ($\approx$1 ano), lendo o beta no
\textbf{último pregão do mês anterior} a cada competência. $\beta>1$ significa que a
ação amplifica o mercado (mais cíclica/arriscada); $\beta<1$, que amortece. A VALE,
mineradora, tem beta tipicamente alto.

\section{Cuidado com \emph{look-ahead} e por que preço \emph{nominal}}

\emph{Look-ahead} é usar informação do futuro para prever o passado — um erro grave.
Garantimos que cada característica olha só para trás (a janela do beta termina na
data de medição; o preço/beta vêm do mês \emph{anterior}). Um detalhe fino: o preço
\emph{ajustado} (adjclose) é recalculado para trás usando dividendos \emph{futuros},
então o nível dele numa data passada ``sabe'' de proventos posteriores. Por isso, na
regressão usamos o \textbf{preço nominal} (o preço de fato negociado, imune a esse
problema); o ajustado serve só para os retornos/beta, que são limpos.

\section{O que os resultados dizem (em palavras)}

O \textbf{nível} do peso de VALE3 é explicado por características \textbf{estruturais}
do fundo: fundos maiores diversificam mais (menos VALE3); fundos com mais cotistas
têm mais; FICs têm menos posição direta. O fluxo do mês não explica o nível (deve
explicar a \emph{variação}). E o ``clima de mercado'' da VALE (preço e risco) desloca
o nível médio de alocação no tempo. Tudo isso é estável e forte ao longo de 6 anos.

\section{Rigor: validação e parsing}

As bases da CVM têm armadilhas (formatos de número com ``R\$'', vírgula decimal,
notação científica; datas em formatos diferentes; linhas duplicadas; mudança de
layout em 2021). Tratamos e \textbf{auditamos} cada uma: 0 duplicados, 0 NAs
indevidos; a captação da SH bate com o Informe Diário da CVM em 99,4--99,8\% ao
centavo nos 6 anos; o preço bate com o fechamento oficial da B3; o beta confere por
três métodos independentes. O recorte de 2016 do painel de 6 anos é idêntico ao
painel de 2016 validado isoladamente.

\section{Passo 2: a dinâmica dos fatores $\theta_t$}

Com os 72 vetores $\theta_t$ em mãos, perguntamos: \emph{como cada coeficiente se
comporta no tempo?} Isso decide o modelo dinâmico.

\textbf{O que achamos.} Os fatores estruturais (intercepto, tamanho, cotistas, FIC)
são \textbf{persistentes mas estacionários}: o coeficiente AR(1) é $\phi \approx
0{,}7$--$0{,}84$ (não $\approx 1$), e um teste tipo Dickey--Fuller rejeita a hipótese
de passeio aleatório. Em palavras: eles \textbf{oscilam em torno de um nível estável
e revertem à média}, em vez de vagar sem rumo. O coeficiente do fluxo é quase ruído
(baixa persistência).

\textbf{Random walk ou reversão à média?} O orientador sugeriu, como ponto de
partida, o passeio aleatório ($\theta_t = \theta_{t-1} + \text{choque}$). Os dados
refinam isso: a dinâmica é de \textbf{reversão à média} (AR(1) com $\phi$ moderado).
Na prática, para prever os fatores estruturais, AR(1) e random walk dão erro
parecido (ambos bons, AR(1) um pouco melhor); o fluxo, por ser ruidoso, é melhor
previsto pela própria média. Ou seja: a posição não ``foge''; ela ancora.

\textbf{Por que então um filtro de Kalman?} As estimativas mensais $\theta_t$ vêm de
regressões com $\sim$130 fundos cada — têm \emph{ruído amostral}. O filtro de Kalman
trata $\theta_t$ como um \emph{estado oculto} que evolui suavemente (AR(1)) e que
``enxergamos'' com ruído a cada mês. Ele combina, de forma ótima, (i) a inércia do
estado e (ii) a informação nova de cada mês, produzindo uma trajetória mais limpa e
uma previsão dos $\theta_{t+1}$ — e, daí, do peso de cada fundo. É o passo seguinte.

\textbf{O produto final} é uma \emph{matriz de erros} (resíduos fundo $\times$ mês),
cujas propriedades (média zero, estabilidade, variância) dizem o quão bem o modelo
captura as posições — exatamente o que o orientador descreveu.

\end{document}
